\documentclass[11pt]{article}
\usepackage[margin=1in]{geometry}
\usepackage{amsmath,amssymb}
\usepackage{booktabs}
\usepackage[hyphens]{url}
\usepackage{hyperref}
\usepackage{newunicodechar}
\newunicodechar{â}{\ensuremath{\hat{a}}}
\newunicodechar{τ}{\ensuremath{\tau}}
\newunicodechar{Γ}{\ensuremath{\Gamma}}
\newunicodechar{θ}{\ensuremath{\theta}}
\newunicodechar{Ω}{\ensuremath{\Omega}}

\setlength{\parskip}{2pt}
\newcommand{\Aoo}{A_{oo}}
\newcommand{\Aob}{A_{ob}}
\newcommand{\Abo}{A_{bo}}
\newcommand{\Abb}{A_{bb}}
\newcommand{\uu}{\mathbf{u}}
\newcommand{\ff}{\mathbf{f}}

\title{\vspace{-2.5em}Element Learning for Spectral--Element Elliptic Solves:\\
       Method and Benchmarking}
\author{Jexpresso}
\date{}

\begin{document}
\maketitle
\vspace{-2.5em}

\section{Governing problem and discretization}
We solve the variable--coefficient Poisson problem on a domain $\Omega\subset\mathbb{R}^2$,
\begin{equation}
  -\nabla\!\cdot\!\big(a(\mathbf{x})\,\nabla u\big) = f \quad\text{in }\Omega,
  \qquad u = g \ \text{on }\partial\Omega ,
  \label{eq:pde}
\end{equation}
with a continuous \emph{spectral--element method} (SEM). The mesh has $N_e$ quadrilateral
elements; each carries a tensor--product nodal basis of degree $k$ on the
Gauss--Lobatto--Legendre (GLL) points, i.e.\ $(k{+}1)^2$ nodes per element. Global
assembly (direct stiffness summation) of the diffusion operator and the mass matrix
yields the sparse linear system
\begin{equation}
  A\,\uu = \mathbf{b}, \qquad \mathbf{b} = M\ff ,
  \label{eq:linsys}
\end{equation}
where $A$ is the SEM Laplacian/diffusion matrix, $M$ the (lumped) mass matrix, and the
Dirichlet data $g$ has been lifted into $\mathbf{b}$. Because \eqref{eq:pde} is
time--independent, \eqref{eq:linsys} is solved \emph{once}.

\paragraph{Baseline: direct SEM solve.} The reference solver factorizes $A$ once
(sparse LU) and back-substitutes: $\uu = A\backslash\mathbf{b}$. The factorization is the
dominant cost; the triangular back-solve is comparatively cheap.

\section{Static condensation}
Partition the degrees of freedom into three groups: the \emph{Dirichlet} nodes
$\Gamma$ (values $g$ known); the element \emph{interior} nodes (the $(k{-}1)^2$ nodes
strictly inside each element, subscript $o$); and the \emph{skeleton} nodes shared
between elements (subscript $b$, the unknown skeleton is denoted $\partial O$ and the
full skeleton $\partial\tau=\partial O\cup\Gamma$). Interior nodes of one element are
coupled only to that element's own skeleton, so the interior unknowns can be eliminated
element by element. Per element $e$, with interior/boundary blocks
$\Aoo^e,\Aob^e,\Abo^e$, define the local \emph{condensation operator} and interior load
\begin{equation}
  T^e = \big(\Aoo^e\big)^{-1}\Aob^e \in \mathbb{R}^{(k-1)^2\times 4k},
  \qquad
  t^e = \big(\Aoo^e\big)^{-1}\ff^{\,e}_o .
  \label{eq:Te}
\end{equation}
Eliminating the interiors produces the condensed \emph{Schur complement} system on the
skeleton unknowns,
\begin{equation}
  S\,\uu_{\partial O} = \hat{\mathbf{b}}_{\partial O},
  \qquad
  S = \Abb - \sum_{e} \Abo^e\,T^e,
  \qquad
  \hat{\mathbf{b}}_{\partial O} = \mathbf{b}_{\partial O} - \sum_e \Abo^e\,t^e
      - S_{\partial O,\Gamma}\,g ,
  \label{eq:schur}
\end{equation}
which is much smaller and better conditioned than $A$. The interior values are then
recovered locally, in parallel over elements,
\begin{equation}
  \uu^{\,e}_o = t^e - T^e\,\uu^{\,e}_b .
  \label{eq:recover}
\end{equation}
Equations \eqref{eq:Te}--\eqref{eq:recover} reproduce the direct solve exactly; the only
per--element work that depends on the geometry/coefficient is forming $T^e$ (one small
dense inverse per element).

\section{Element learning}
The key observation is that $T^e$ in \eqref{eq:Te} is a smooth function of the local
element geometry and diffusivity alone. \emph{Element learning} replaces the per-element
inverse by a trained neural network surrogate
\begin{equation}
  T^e \;\approx\; \mathcal{N}_\theta\!\big(\hat a^{\,e}\big),
  \qquad
  \hat a^{\,e}\in\mathbb{R}^{(k+1)^2},
  \label{eq:surrogate}
\end{equation}
where the feature $\hat a^{\,e}$ is the diffusivity $a(\mathbf{x})$ sampled at the
element's $(k{+}1)^2$ nodes (for constant $a$ a single scalar broadcast to all nodes;
for the non-constant case a geometry-induced anisotropic $\hat a$). The network maps
$(k{+}1)^2$ inputs to the $4k(k{-}1)^2$ entries of $\mathrm{vec}(T^e)$. Only the
\emph{operator} $T^e$ is learned; the interior load $t^e$ in \eqref{eq:Te} is still formed
exactly from the stored interior block and $\ff^e_o$, so the source term is never
approximated. With $\mathcal{N}_\theta$ in hand, \eqref{eq:schur}--\eqref{eq:recover}
proceed as before, but each $T^e$ is a forward network evaluation rather than a dense
inverse.

\paragraph{Training data.} The map \eqref{eq:surrogate} is learned \emph{offline} on a
single reference element. Sampling runs the SEM assembly on a $1\times1$ mesh: for many
random diffusivities $\hat a$ it computes the exact $T^e$ from \eqref{eq:Te} and stores the
pairs $(\hat a,\,T^e)$ as \texttt{input\_tensor.csv} / \texttt{output\_tensor.csv}. The
network (a fully connected MLP / small CNN) is trained with input and output
standardization, then exported to ONNX so the Julia solver reads the raw feature and
returns the raw $T^e$.

\section{Offline--online pipeline}
A single launcher automates the three stages end to end:
\begin{enumerate}\itemsep2pt
  \item \textbf{Sample} --- run SEM in sampling mode (\texttt{lEL\_Sample=true}) on the
        $1\times1$ mesh $\Rightarrow$ tagged training CSVs (one pair per case).
  \item \textbf{Train} --- the Python trainer turns the CSVs into an ONNX model
        $\mathcal{N}_\theta$.
  \item \textbf{Infer} --- re-run SEM in inference mode (\texttt{lEL\_Sample=false}) on the
        real multi-element mesh, loading $\mathcal{N}_\theta$ to build every $T^e$ via
        \eqref{eq:surrogate}, then solving \eqref{eq:schur}--\eqref{eq:recover}.
\end{enumerate}
The learned operator is mesh-size independent: trained once on the reference element, it
is reused for every element of an arbitrarily large mesh.

\section{Cost model and benchmarking}
Each method is timed at two levels (Table~\ref{tab:cost}); speedup is
$t_{\mathrm{SEM}}/t_{\mathrm{EL}}$. Each level is a separate \texttt{@btime}
measurement of a specific callable, so the two rows differ precisely in
\emph{which work is inside the timed call} versus done once beforehand.

\paragraph{This solve (from scratch).} The full cost of the one solve that runs.
\emph{SEM} times \verb|A\b|, which for a sparse $A$ computes a fresh LU
factorization and then back-substitutes on every sample. \emph{EL} times the whole
condensation solve \verb|elementLearning_Axb!|: extract the per-element blocks
$\{\Aoo^e,\Aob^e,\Abo^e,\Abb\}$ from the global $A$, evaluate the surrogate
$\{T^e\}$, assemble the Schur complement $S$ \eqref{eq:schur}, solve the skeleton
system, and recover the interiors \eqref{eq:recover}. Both include all of their
$A$-dependent work, so this is the honest single-solve comparison.

\paragraph{Amortized (setup reused).} The per-solve cost once each method's
\emph{entire} $A$-dependent setup has been done once and cached --- a true
apples-to-apples comparison in which each method simply back-substitutes its own
cached factorization for a new right-hand side. \emph{SEM} pre-computes
$F=\mathrm{lu}(A)$ and times only \verb|F\b| (the full system). \emph{EL}, having
already assembled the surrogate operators $\{T^e\}$ and the Schur complement $S$
during the solve above, pre-computes $F_S=\mathrm{lu}(S_{\partial O,\partial O})$
and times only $F_S\backslash\mathbf{b}$ --- the back-substitution of the
\emph{condensed skeleton} system \eqref{eq:schur}, which is much smaller than $A$
(only the inter-element skeleton unknowns). This isolates the core linear-solve
cost: the direct method back-solves the full $N\times N$ system, element learning
back-solves the reduced skeleton system.

\begin{table}[h]
\centering
\small
\begin{tabular}{@{}lll@{}}
\toprule
Level (timed call) & Direct SEM & Element learning \\
\midrule
\textbf{This solve} \verb|A\b| / \verb|Axb!| & factorize $A$ + back-solve & extract blocks + surrogate + assemble $S$ + skeleton solve + recover \\[2pt]
\textbf{Amortized} \verb|F\b| / \verb|Fs\b| & back-solve full system, $F=\mathrm{lu}(A)$ cached & back-solve skeleton system, $F_S=\mathrm{lu}(S)$ cached \\
\bottomrule
\end{tabular}
\caption{What is inside each timed call. In the \textbf{Amortized} row both methods
reuse their full $A$-dependent setup (SEM's factorization of $A$; EL's block
extraction, surrogate $\{T^e\}$, Schur assembly, and factorization of $S$) and time
only the back-substitution for a new right-hand side --- SEM on the full $N\times N$
system, EL on the smaller condensed skeleton system. Element learning's advantage in
this regime comes from that size reduction. (The full EL per-RHS map also includes
forming the condensed load and recovering the interiors, both $O(N_e)$ cheap,
element-local, and parallel; only the dominant reduced back-solve is timed here.)}
\label{tab:cost}
\end{table}

\noindent The global matrix $A$ is assembled once before either solve and shared, so it
is counted for neither. All solver times are reported as \texttt{@btime} minima
(BenchmarkTools: many samples, compilation excluded) so they are repeatable rather than
a noisy single shot. Accuracy is measured in the mass-weighted $L_2$ norm
$\|e\|_M=(\sum_i M_i e_i^2)^{1/2}$ (absolute and relative) and $L_\infty$, comparing the
inferred solution against (i) the direct SEM solution, and (ii) the exact/manufactured
solution $q_e$ when the case provides one. The inference run additionally writes the SEM
and exact fields and their differences from the inferred solution into the VTU output for
side-by-side visualization.

\end{document}
